% \begin{abstract}
% We show that training on image editing tasks can systematically improve image understanding \amil{maybe you can see on tasks such as X,Y,Z. I'm not sure exactly what does and doesn't count as ``image understanding``}.
% % While it is widely known that understanding helps generation, the reverse of whether and how generation helps understanding remains poorly understood.
% We first construct tasks that share the same underlying visual relation but can be trained either as image-to-image editing (I2I) or image-to-text prediction (I2T). This allows us to directly test whether I2I supervision transfers to I2T understanding. Across these controlled settings, we find that I2I training consistently improves I2T performance.
% We then ask which generation skills transfer to which understanding capabilities. We organize image-editing data using established visual task taxonomies, and regroup existing vision-language benchmarks into a capability-based evaluation taxonomy.
% Through training sweeps, we identify a recipe that improves transfer from I2I supervision to I2T evaluation.
% Applying this recipe reveals a structured transfer map between generation and understanding: for example, depth prediction most strongly improves depth and spatial reasoning, while keypoint-based editing emerges as a broadly transferable task that benefits multiple capabilities.
% Finally, we use these insights to train a unified model \amil{to potentially boost impact you can name your trained model and the size} that achieves stronger understanding ability by only adding more I2I data as supervision. These results suggest that image editing is not merely a generative objective, but also a useful source of supervision for visual understanding.
% \end{abstract}

\begin{abstract}
% Training a mdoel to \textit{generate} visual content implicitly teaches it about geometry, space, objectness, and other perceptual capabilities; yet, its benefits for visual \textit{understanding} remains unclear. We ask: how does visual generation supervision improve visual understanding?
% \xd{note from xd: all changes are directly colored as blue.}
Training a model to \textit{generate} visual content
can encourage it to learn rich perceptual capabilities related to
geometry, spatial relationships, and objectness;
yet, its benefits for visual \textit{understanding} remain unclear.
We ask: when and how does visual generation supervision improve visual understanding?
%% controlled study
% Using controlled pairs of image-to-image (I2I) generation and image-to-text (I2T) understanding tasks that are two variants of the same problem, we find that I2T performance improves from I2I supervision under an appropriate curriculum of first pre-training I2I, with \textit{gains increasing as more I2I data is added}.
We study controlled pairs of image-to-image (I2I) generation and
image-to-text (I2T) understanding tasks that express the same underlying problem in different output modalities. 
We find that under the correct recipe, I2I training improves downstream I2T performance, with \textit{larger gains as the amount of I2I training data increases}.
%% taskonomy
% We next ask which generation tasks benefit which understanding tasks.
% To study transfer beyond paired tasks, we introduce \methodname, a unified taxonomy spanning 19 I2I tasks and 25 T2I understanding capabilities. The resulting transfer map reveals selective benefits: some follow intuitive correspondences, such as depth supervision improving metric 3D reasoning, object pointing improving counting, and Jigsaw improving 2D ordering; meanwhile, some are less intuitive, such as colorization improving visual correspondence, 2D edges improving depth understanding.
We next ask which generation tasks benefit which understanding capabilities. To study transfer beyond paired tasks, we introduce \methodname, a unified taxonomy spanning 19 I2I generation tasks and 25 I2T understanding capabilities. The resulting transfer map reveals selective, task-dependent benefits. Some follow intuitive correspondences, \textit{e.g.,} depth prediction improving metric 3D reasoning, object pointing improving counting, and jigsaw reconstruction improving 2D ordering. 
Interestingly, we also uncover surprising connections: 2.5D segmentation improving category recognition and Z-depth prediction improving localization. 
% We find that transfer can follow intuitive correspondences, such as depth improving metric 3D reasoning and object pointing improving counting, but can also extend to less directly related capabilities, such as keypoints improving lighting understanding.
% The same generation task can also benefit some targets while interfering with others.
%% gradient
% To better understand these transfer patterns, we mechanistically analyze the gradients of I2I and I2T task pairs, to find that capabilities with higher gradient alignment at these layers also exhibit larger downstream transfer gains.
To probe these patterns, we analyze gradient alignment between generation and understanding tasks and find that stronger alignment is associated with larger downstream transfer gains.
% Together, our results show that generation supervision can substantially improve visual understanding with a specific training strategy, and that different generation tasks benefit different understanding capabilities. 
Together, our results highlight visual generation as a rich source of supervision for visual understanding and provide a roadmap for unlocking its benefits through the right training curriculum and task selection. Project page: \url{https://omni-taskonomy.github.io/}.
\end{abstract}
% start with "taxonomy" and "bridging" make the abs more exciting; start with stating the position

% we need to construct briding tasks and taxonomy

% then mention controlled settings
% we did I2I helps I2T


% gradients analyisis; manifold 训两个obj, 是不是两个obj让gradient往一个方向; 否则迭代能力orthogonal; I2I task可以模拟I2T sft时候的traj, 所以I2I的optimization方向更加合理; 注意interp的时候不要update weights; 

% when does image generation helps image understanding,?when it doesn't?

% when does image generation helps image understanding,?

% xxx image editing helps image understanding

% when does image editing helps image understanding?

% when does image generation helps image understanding and why?
